% ============================================================================
% FARZULLA RESEARCH PREPRINT TEMPLATE
% ============================================================================
% Custom academic preprint template with branding, version control, and logos
% Author: Murad Farzulla
% Organization: Farzulla Research
% Template Version: 1.0.0
% Last Updated: November 2025
%
% QUICK START - CUSTOMIZE THESE PLACEHOLDERS PER PAPER:
% 1. [Your Paper Title Here] - Full title
% 2. [Your Subtitle Here] - Subtitle (italic)
% 3. [Short Title Here] - Running header (2-4 words, concise)
% 4. [Month Year] - Publication date (e.g., "November 2025")
% 5. [keyword1, keyword2, keyword3] - PDF metadata keywords
%
% Author info already filled in (Murad Farzulla, ORCID, Farzulla Research, email)
%
% FEATURES IMPLEMENTED:
% - Corner logos (FR left, Zenodo right) aligned with preprint banner
% - Burgundy section headings (keeps sizes unchanged)
% - Running headers: Author (left) | Short Title (right)
% - Academic affiliation footnotes (superscript on institution, clickable)
% - URL line breaking (fixes long URLs in bibliography)
% - Hyperlinked citations (clickable in-text → bibliography)
% - First page style (no header, just logos + footers)
% ============================================================================

\documentclass[11pt,twocolumn]{article}

% ============================================================================
% PACKAGE IMPORTS
% ============================================================================

% Page geometry and layout
\usepackage[a4paper, margin=1in]{geometry}
\usepackage{setspace}

% Fonts and typography
\usepackage[T1]{fontenc}
\usepackage{lmodern}

% Graphics and figures
\usepackage{graphicx}
\usepackage{float}
\graphicspath{{arxiv-submission/figures/}}

% Tables
\usepackage{booktabs}
\usepackage{array}
\usepackage{longtable}
\usepackage{multirow}

% Lists
\usepackage{enumitem}

% Math packages
\usepackage{amsmath}
\usepackage{amssymb}
\usepackage{amsthm}

% Bibliography
\usepackage[round,authoryear]{natbib}
\bibliographystyle{abbrvnat}

% Colors
\usepackage{xcolor}
\definecolor{farzullaburgundy}{RGB}{128,0,32}      % Farzulla Research brand color
\definecolor{zenodoblue}{RGB}{0,123,255}            % Zenodo blue
\definecolor{orcidgreen}{RGB}{166,206,57}           % ORCID green

% Colored boxes
\usepackage{tcolorbox}

% TikZ for custom graphics
\usepackage{tikz}
\usepackage{scalerel}

% Hyperlinks and PDF metadata
\usepackage{url}  % Better URL formatting and line breaking
\usepackage[colorlinks=true,
            linkcolor=farzullaburgundy,
            citecolor=farzullaburgundy,
            urlcolor=farzullaburgundy,
            breaklinks=true,
            pdftitle={Testing Narrative-Market Alignment in Cryptocurrency: A Methodological Framework},
            pdfauthor={Murad Farzulla},
            pdfkeywords={cryptocurrency, tensor decomposition, natural language processing, market microstructure, whitepaper analysis, factor analysis, Procrustes rotation, Tucker congruence}]{hyperref}

% Allow URLs to break at hyphens and slashes
\def\UrlBreaks{\do\/\do-\do_}
\expandafter\def\expandafter\UrlBreaks\expandafter{\UrlBreaks\do\a\do\b\do\c\do\d\do\e\do\f\do\g\do\h\do\i\do\j\do\k\do\l\do\m\do\n\do\o\do\p\do\q\do\r\do\s\do\t\do\u\do\v\do\w\do\x\do\y\do\z}

% Section spacing and formatting
\usepackage{titlesec}
\titlespacing*{\section}{0pt}{1.5ex plus 0.5ex minus 0.2ex}{1ex plus 0.2ex}
\titlespacing*{\subsection}{0pt}{1.2ex plus 0.4ex minus 0.2ex}{0.8ex plus 0.2ex}

% Burgundy section headings (justified for traditional block look, scaled down)
\titleformat{\section}{\normalfont\large\bfseries\color{farzullaburgundy}}{\thesection}{0.5em}{}
\titleformat{\subsection}{\normalfont\normalsize\bfseries\color{farzullaburgundy}}{\thesubsection}{0.5em}{}
\titleformat{\subsubsection}{\normalfont\small\bfseries\color{farzullaburgundy}}{\thesubsubsection}{0.5em}{}

% Footer customization
\usepackage{fancyhdr}

% Absolute positioning for corner logos
\usepackage{eso-pic}

% ============================================================================
% CUSTOM LOGO COMMANDS
% ============================================================================

% ORCID icon (clickable green circle with iD mark) - adjustable size
\newcommand{\orcidicon}{\scalerel*{
    
\begin{tikzpicture}[x=3ex,y=3ex]
    \draw[fill=orcidgreen] (0,0) circle (0.5);
    \draw[white,line width=0.08ex] (0,0.15) -- (0,0.3);
    \draw[white,line width=0.08ex] (-0.15,0) -- (-0.3,0);
    \end{tikzpicture}
}{\textrm{I}}}

% Farzulla Research logo (actual logo image - PDF for vector scaling)
\newcommand{\farzullalogo}{%
    \href{https://farzulla.org}{%
        \includegraphics[height=4ex]{farzulla-logo.pdf}%
    }%
}

% Zenodo logo (actual logo image - PDF for vector scaling)
\newcommand{\zenodologo}{%
    \href{https://zenodo.org}{%
        \includegraphics[height=4ex]{zenodo-logo.pdf}%
    }%
}

% ============================================================================
% VERSION AND PREPRINT BANNER
% ============================================================================

% Version number command (CUSTOMIZE - update for each version/paper)
\newcommand{\paperver}{2.0.1}
\newcommand{\paperdate}{December 2025}

% Corner logos removed - using header/footer text instead

% ============================================================================
% METADATA BOX (appears after abstract on first page)
% ============================================================================

\newcommand{\metadatabox}[1]{%
\begin{tcolorbox}[
    colback=gray!5,
    colframe=farzullaburgundy,
    title=Publication Metadata,
    fonttitle=\bfseries,
    coltitle=white,
    colbacktitle=farzullaburgundy,
    width=\columnwidth,
    arc=2mm,
    boxrule=0.8pt
]
\small
\textbf{DOI:} \href{https://doi.org/#1}{\texttt{#1}}\\
\textbf{Version:} \paperver\\
\textbf{Date:} \paperdate\\
\textbf{License:} CC-BY-4.0
\end{tcolorbox}
}

% ============================================================================
% HEADER AND FOOTER CONFIGURATION
% ============================================================================

% Short title command (CUSTOMIZE THIS - used in running header)
% Keep it concise (2-4 words max)
\newcommand{\shorttitle}{Narrative-Market Alignment Framework}

\pagestyle{fancy}
\fancyhf{} % Clear default

% Running headers (for main body pages)
\fancyhead[L]{\small\href{https://farzulla.org}{farzulla.org}}  % Left: Website
\fancyhead[C]{}                                                  % Center: Empty
\fancyhead[R]{\small\href{https://doi.org/10.5281/zenodo.17772652}{DOI: 10.5281/zenodo.17772652}}  % Right: DOI

% Footers
\fancyfoot[C]{\small\thepage}
\fancyfoot[L]{\small Murad Farzulla}
\fancyfoot[R]{\small v\paperver~| \paperdate}

\renewcommand{\headrulewidth}{0.4pt}
\renewcommand{\footrulewidth}{0.4pt}

% First page style (no headers, just footers)
\fancypagestyle{firstpage}{
  \fancyhf{}
  \fancyfoot[C]{\small\thepage}
  \fancyfoot[L]{\small Murad Farzulla}
  \fancyfoot[R]{\small v\paperver~| \paperdate}
  \renewcommand{\headrulewidth}{0pt}
  \renewcommand{\footrulewidth}{0.4pt}
}

% ============================================================================
% DOCUMENT METADATA
% ============================================================================

\title{\vspace{-1em}\textbf{Tensor Structure in Cryptocurrency Microstructure: Low-Rank CP Decomposition from Binance OHLCV}\\[0.5em]
\large \textit{Single-Venue Proof-of-Concept}\vspace{-0.5em}}

\author{
    \textbf{Murad Farzulla} \href{https://orcid.org/0009-0002-7164-8704}{\orcidicon}\\
    \textit{Farzulla Research} \farzullalogo\\[0.3em]
    \small \paperdate
}

\date{}

% ============================================================================
% DOCUMENT BODY
% ============================================================================

\begin{document}

% Single-column frontmatter
\onecolumn
\setstretch{1.2}

% Apply first page style (no header)
\thispagestyle{firstpage}

% Preprint status centered above title
\begin{center}
{\small\color{farzullaburgundy}\textbf{PREPRINT v\paperver} | \color{gray}Not peer-reviewed}
\end{center}
\vspace{0.5em}

% Manual title construction
\begin{center}
{\Large\bfseries Testing Narrative-Market Alignment in Cryptocurrency:}\\[0.3em]
{\Large\bfseries A Methodological Framework}\\[0.5em]
{\large\itshape Testing Alignment Between Functional Taxonomies and Market Dynamics}\\[1em]
{\bfseries Murad Farzulla}\textsuperscript{1} \href{https://orcid.org/0009-0002-7164-8704}{\orcidicon\ \texttt{0009-0002-7164-8704}}\\[0.5em]
{\small\itshape \textsuperscript{1}\href{https://farzulla.org}{Farzulla Research}}\\[0.3em]
{\small \paperdate}\\[0.5em]
{\small Correspondence: \href{mailto:murad@farzulla.org}{murad@farzulla.org}}
\end{center}

% Abstract
\begin{abstract}
\noindent \textit{Version 2.0.1 (December 2025) - Proposes a methodological framework for testing whether functional claims in cryptocurrency whitepapers align with market factor structures.}

\textbf{Research Question:} Do the functional purposes articulated in cryptocurrency whitepapers (store of value, smart contracts, infrastructure, etc.) correspond to how markets actually price these assets?

\textbf{Methodological Contribution:} We propose a three-stage framework combining natural language processing with market characterization to test narrative-market alignment. Stage 1: Extract functional profiles from whitepapers using zero-shot classification (BART-MNLI) against a 10-category functional taxonomy. Stage 2: Construct market profiles from trading data (returns, volatility, volume metrics). Stage 3: Measure alignment between claims and market spaces using Procrustes rotation and Tucker's congruence coefficient ($\phi$).

\textbf{Illustrative Application:} We demonstrate the framework on a pilot sample of 8 cryptocurrencies with 365 days of hourly Binance data. The pilot yields $\phi = 0.719$ (95\% CI: [0.623, 0.953]), interpreted as ``some similarity.'' However, we emphasize that this N=8 sample is insufficient for robust empirical inference---our contribution is the \textit{methodology}, not the point estimates. The wide confidence interval and limited sample preclude strong claims about alignment magnitude.

\textbf{Key Insight:} Bitcoin's whitepaper emphasizes ``peer-to-peer electronic cash'' while markets price BTC as ``digital gold''---illustrating the core problem. Whitepapers capture \textit{static founding mandates}, while market profiles capture \textit{evolved market utility}. The gap between them represents the empirically testable alignment question this framework addresses.

\textbf{Scope:} This paper establishes methodological foundations for narrative-market alignment testing. Future work with expanded samples (N $\geq$ 30) is required for robust empirical conclusions.
\end{abstract}

\vspace{0.5em}
\noindent \textbf{Keywords:} Cryptocurrency, tensor decomposition, natural language processing, market microstructure, whitepaper analysis, factor analysis, Procrustes rotation, Tucker congruence

\vspace{0.5em}
\noindent \textbf{JEL Classification:} C38 (Classification Methods), G12 (Asset Pricing), C45 (Neural Networks), G14 (Information and Market Efficiency)

% Metadata box with DOI
\metadatabox{10.5281/zenodo.17772652}

\vspace{0.5em}
\noindent \textbf{Code \& Data:} All analysis code, whitepapers, and results available at \url{https://github.com/studiofarzulla/tensor-defi}

\section*{Research Context}

This work forms part of the Adversarial Systems Research program, which investigates stability, alignment, and friction dynamics in complex systems where competing interests generate structural conflict. The program examines how agents with divergent preferences interact within institutional constraints across multiple domains: political governance, financial markets (cryptocurrency volatility and regulatory responses), human cognitive development (trauma as maladaptive learning from adversarial training environments), and artificial intelligence alignment (multi-agent systems with competing objectives).

In cryptocurrency markets, adversarial dynamics manifest as the tension between \textit{founding mandates} (what projects claim to do in whitepapers) and \textit{evolved market utility} (how participants actually price assets over time). Whitepapers serve as static founding documents that articulate functional purposes---store of value, smart contracts, infrastructure, governance---yet markets may price assets according to criteria that \textit{evolved} post-launch: liquidity, speculation, network effects, ecosystem development, or regulatory exposure.

\textbf{Contribution:} This paper proposes a methodological framework for testing narrative-market alignment. We extract functional taxonomies from whitepapers using NLP and construct market factor structures via tensor decomposition, then measure alignment using Procrustes rotation and Tucker's congruence coefficient. The pilot application (N=8 entities) yields $\phi = 0.719$, but we emphasize the \textit{method}, not the estimate---the small sample precludes strong empirical claims.

\textbf{Key Framing:} We interpret divergence not as ``marketing vs. reality'' (which implies deception) but as \textit{static mandate vs. evolved utility}. Bitcoin's whitepaper describes ``peer-to-peer electronic cash,'' but the asset's \textit{utility} evolved toward ``digital gold'' as market participants found new uses. The gap between founding mandate and current utility is precisely what this framework measures.

% Table of contents on new page
\clearpage
\tableofcontents

% ====================
% SWITCH TO TWO-COLUMN LAYOUT
% ====================
\clearpage
\twocolumn
\setstretch{1.5}  % Re-apply line spacing after column switch

% ====================
% MAIN CONTENT
% ====================

\section{Introduction}
\label{sec:introduction}

Cryptocurrency whitepapers serve as founding documents that articulate what a project claims to do: store of value, medium of exchange, smart contract platform, decentralized infrastructure, privacy-preserving transactions, governance mechanisms, oracle services, identity systems, gaming/metaverse applications, or stablecoins. These functional claims shape investor expectations, regulatory classification, and market positioning.

But do markets actually price assets according to these stated functions? Or do market dynamics reflect entirely different factors---liquidity, speculation, network effects, regulatory exposure---that diverge from whitepaper narratives?

\subsection{The Research Gap}

Extensive work examines either textual analysis of crypto documentation \citep{Liu2021Risks} or market microstructure dynamics \citep{Almeida2024Cryptocurrency}, but no prior research systematically tests the \textit{alignment} between these domains. This gap matters for several reasons:

\begin{itemize}[leftmargin=*]
    \item \textbf{Market efficiency}: If markets price assets according to stated function, whitepapers contain decision-relevant information. If not, stated purpose may be ``cheap talk.''
    \item \textbf{Due diligence}: Detecting narrative-market misalignment could identify assets where marketing diverges from fundamentals.
    \item \textbf{Regulatory classification}: Functional claims inform whether tokens are securities, commodities, or currencies---alignment testing provides empirical grounding.
\end{itemize}

\subsection{Our Approach}

We combine natural language processing with market characterization to construct and align two representations of the same entities:

\textbf{Claims Space:} Extract functional profiles from whitepapers using zero-shot classification (BART-MNLI) against a 10-category taxonomy. Each entity becomes a vector in $\mathbb{R}^{10}$ representing its claimed functional distribution.

\textbf{Market Space:} Compute market profiles from hourly OHLCV data (returns, volatility, Sharpe ratio, drawdown, volume metrics). Each entity becomes a vector in $\mathbb{R}^{7}$ representing market-revealed characteristics.

\textbf{Alignment Testing:} Apply Procrustes rotation to align the claims and market spaces, then compute Tucker's congruence coefficient ($\phi$) to measure similarity.

\subsection{Core Hypothesis}

\textbf{H$_0$:} Whitepaper functional profiles are orthogonal to market profiles ($\phi < 0.65$)

\textbf{H$_1$:} Whitepaper functional profiles show significant alignment with market profiles ($\phi \geq 0.65$)

A threshold of $\phi = 0.65$ corresponds to ``some similarity'' under standard Tucker congruence interpretation \citep{Lorenzo-Seva2006Tucker}.

\subsection{Contributions}

This paper makes the following contributions:

\begin{enumerate}[leftmargin=*]
    \item \textbf{Methodological framework}: First systematic framework for testing narrative-market alignment in cryptocurrency markets, combining NLP and tensor methods
    \item \textbf{Functional taxonomy}: 10-category classification scheme for cryptocurrency whitepaper claims, applicable to broader crypto text analysis
    \item \textbf{Alignment methodology}: Procrustes rotation + Tucker's congruence for comparing heterogeneous factor spaces (claims vs. market)
    \item \textbf{Illustrative pilot}: Demonstration on N=8 entities showing framework feasibility (note: pilot sample insufficient for robust empirical inference)
    \item \textbf{Conceptual framing}: ``Static mandate vs. evolved utility'' interpretation that distinguishes functional evolution from misrepresentation
\end{enumerate}

\subsection{Organization}

Section~\ref{sec:related} reviews related work on NLP in finance, tensor methods, and alignment testing. Section~\ref{sec:methodology} details the three-stage pipeline: NLP extraction, tensor decomposition, and alignment testing. Section~\ref{sec:results} presents alignment results with bootstrap confidence intervals. Section~\ref{sec:discussion} interprets findings and discusses implications. Section~\ref{sec:conclusion} concludes.

\section{Related Work}
\label{sec:related}

\subsection{Market Microstructure Theory}

Traditional market microstructure theory \citep{OHara1995Market,Hasbrouck2007Empirical} focuses on the mechanisms through which assets are traded and prices are formed. Classic models emphasize bid-ask spreads, order flow dynamics, and information asymmetry. \citet{Easley2012Flow} developed VPIN (Volume-Synchronized Probability of Informed Trading) to measure order flow toxicity, which has been applied to cryptocurrency flash crashes \citep{Easley2024CrossAsset}. However, these frameworks were developed for centralized exchanges with unified order books, and their extension to decentralized, multi-venue cryptocurrency markets remains underexplored.

\subsection{Dimensionality Reduction in Finance}

Principal Component Analysis (PCA) has been the dominant approach for factor extraction in financial markets since the early work of \citet{Cont2001Empirical} on stylized facts of asset returns. \citet{Fan2013Large} extended PCA to high-dimensional settings via covariance thresholding, enabling factor-based risk models when the number of assets exceeds the number of observations. Alternatives to PCA include Independent Component Analysis \citep{Back1997First} and deep autoencoders \citep{Heaton2017Deep}, though these remain matrix-based methods that flatten multi-dimensional data.

Recent cryptocurrency factor studies have identified asset-specific drivers. \citet{Liu2021Risks} established that network growth and momentum factors explain 64\% of cross-sectional cryptocurrency returns, while \citet{Fieberg2023CTREND} proposed a trend factor (CTREND) capturing 52\% of return variation. Industry research by Sparkline Capital highlights intangible value (developer activity metrics) as the strongest factor with a Sharpe ratio of 0.89 \citep{Sparkline2025IntangibleValue}. However, all these approaches employ matrix factorization, treating venue, asset, and temporal dimensions separately rather than jointly.

\subsection{Tensor Decomposition Methods}

Tensor decomposition is well-established in signal processing and machine learning \citep{Kolda2009Tensor,Cichocki2015Tensor}, but financial applications remain scarce. \citet{Han2022CP} developed a CP factor model for dynamic tensor time series with theoretical convergence guarantees, while \citet{Chang2023Modelling} proposed a one-pass CP estimation via generalized eigenanalysis, avoiding iterative optimization. \citet{Babii2023Tensor} proved that tensor PCA achieves optimal convergence rates for financial factor models, with estimation error scaling as $O(\sqrt{\sum_k p_k / n})$ compared to $O(\sqrt{\prod_k p_k / n})$ for matrix PCA--a dramatic improvement when dimensions are large.

Empirical applications demonstrate tensor methods' superiority. \citet{Lettau2024ThreeD} applied 3D-PCA to sorted equity portfolios and achieved 65--76\% explained variance versus 50\% for traditional PCA. \citet{Mo2025ACT} proposed a tensor completion framework for missing financial data, achieving 40--74\% lower imputation error than matrix methods. \citet{Battaglino2018Practical} developed randomized CP-ALS, enabling 5--10x computational speedup with minimal accuracy loss--critical for high-frequency financial data.

Despite these advances, no prior work applies tensor decomposition to multi-asset cryptocurrency market microstructure. The closest predecessor is \citet{Watorek2023Decomposing}, who used matrix PCA to decompose single-asset high-frequency Bitcoin dynamics, manually identifying three daily activity phases (Asian, European, US sessions). Our tensor framework extends this by jointly modeling multiple assets and time--preserving multi-dimensional structure that single-asset analysis destroys. This proof-of-concept establishes the methodological foundation; future versions will expand to multi-venue analysis to capture cross-exchange arbitrage dynamics.

\subsection{Cryptocurrency Market Microstructure}

\citet{Almeida2024Cryptocurrency} systematically reviewed 138 cryptocurrency microstructure papers (2016--2023) and identified a critical research gap: \textbf{no studies apply high-dimensional joint modeling across exchanges}. Existing work either analyzes single venues \citep{Bozzetto2023Microstructure} or performs pairwise comparisons \citep{Wang2024CrossExchange}.

\citet{Bianchi2022Trading} documented substantial heterogeneity in liquidity provision across 1,000+ cryptocurrency pairs and exchanges, with liquidity premia 3x larger for illiquid pairs. \citet{Farzulla2025Cryptocurrency} documented infrastructure dominance over regulatory uncertainty in cryptocurrency volatility (2.385\% vs 0.419\%, $p = 0.0008$, Cohen's $d = 2.753$), with 5$\times$ sensitivity amplification during crisis periods and substantial cross-sectional heterogeneity ranging from 1.19\% (BTC) to 3.37\% (ADA). \citet{Liu2025LiquidityCommonality} found significant liquidity commonality during COVID-19, with market-wide shocks affecting multiple assets simultaneously--evidence of low-rank factor structure. \citet{DeFreitas2025HighFreq} showed that Bitcoin perpetual futures lead spot markets by 0.5--2 seconds, with lead times increasing during volatility--suggesting time-varying cross-market dynamics.

\citet{Wang2024CrossExchange} modeled cross-exchange risk via dynamic correlation networks, finding that network density spikes during crises (Luna collapse, FTX failure). However, their approach requires pairwise correlation matrices, losing higher-order interactions. \citet{Easley2024CrossAsset} demonstrated cross-asset microstructure spillovers using machine learning on lagged metrics, but their pairwise analysis cannot capture simultaneous three-way interactions (BTC-ETH-SOL spillovers).

\subsection{Our Contribution}

This paper contributes a novel framework for testing \textit{narrative-market alignment}---whether whitepaper claims correspond to market factor structures. Our approach differs from prior work in three ways:

\begin{enumerate}[leftmargin=*]
    \item \textbf{Cross-domain alignment}: We bridge NLP (whitepaper text) and quantitative finance (market tensors), testing whether stated function predicts market positioning.
    \item \textbf{Tensor-based factor extraction}: We preserve the multi-dimensional structure (asset $\times$ time $\times$ feature) rather than flattening to matrices, extracting interpretable market factors via CP decomposition.
    \item \textbf{Formal alignment testing}: We apply Procrustes rotation and Tucker's congruence to measure correspondence between heterogeneous factor spaces (claims vs. market), providing a rigorous statistical framework.
\end{enumerate}

Building on the theoretical foundations of \citet{Babii2023Tensor} and the empirical success of \citet{Lettau2024ThreeD}, we extract market factors via tensor decomposition and test their alignment with NLP-derived functional profiles---a novel application of tensor methods to narrative-market correspondence.

\section{Methodology}
\label{sec:methodology}

Our methodology comprises three stages: (1) NLP pipeline for extracting functional profiles from whitepapers, (2) market profile construction from trading data, and (3) alignment testing via Procrustes rotation and Tucker's congruence.

\subsection{Stage 1: NLP Pipeline}

\subsubsection{Whitepaper Collection}

We collect whitepapers for the top 8 cryptocurrencies by market capitalization with available documentation:

\begin{itemize}[leftmargin=*]
    \item \textbf{Sources}: AllCryptoWhitepapers.com, official project documentation, GitHub repositories
    \item \textbf{Universe}: BTC, ETH, SOL, AVAX, DOT, FIL, LINK, ALGO
    \item \textbf{Format}: PDF documents extracted via PyMuPDF
\end{itemize}

\subsubsection{Text Extraction and Chunking}

Documents are processed into semantic chunks:
\begin{itemize}[leftmargin=*]
    \item Maximum 512 tokens per chunk
    \item Paragraph boundaries preserved
    \item Headers, footers, and page numbers removed
\end{itemize}

\subsubsection{Zero-Shot Classification}

We apply zero-shot classification using \texttt{facebook/bart-large-mnli} against a 10-category functional taxonomy:

\begin{table*}[htbp]
\footnotesize
\caption{Functional Taxonomy for Whitepaper Classification}
\label{tbl:taxonomy}
\begin{tabular}{@{}p{3.2cm}p{4.5cm}p{7cm}@{}}
\toprule
\textbf{Category} & \textbf{Description} & \textbf{Example Claims} \\
\midrule
store\_of\_value & Digital store of value, savings & ``digital scarcity'', ``deflationary'', ``inflation hedge'' \\
medium\_of\_exchange & Payment system, transactions & ``micropayments'', ``peer-to-peer cash'', ``remittances'' \\
smart\_contracts & Programmable contracts, dApps & ``programmable money'', ``composability'', ``DeFi'' \\
infrastructure & Blockchain infra, scaling & ``layer 1'', ``sharding'', ``rollups'', ``interoperability'' \\
privacy & Privacy-preserving transactions & ``zero-knowledge'', ``confidential'', ``anonymous'' \\
governance & Decentralized governance, DAOs & ``on-chain governance'', ``staking'', ``voting'' \\
data\_oracle & External data feeds, oracles & ``price feeds'', ``external data'', ``off-chain'' \\
identity & Digital identity, authentication & ``DID'', ``verifiable credentials'', ``self-sovereign'' \\
gaming\_metaverse & Gaming, virtual worlds & ``play-to-earn'', ``metaverse'', ``NFT'' \\
stablecoin & Price-stable assets & ``pegged'', ``algorithmic stable'', ``collateralized'' \\
\bottomrule
\end{tabular}
\end{table*}

For each chunk, we obtain classification scores for all categories (multi-label). Scores above threshold 0.3 are aggregated and normalized to produce entity-level functional profiles.

\subsubsection{Claims Matrix Construction}

The output is a claims matrix $\mathbf{X} \in \mathbb{R}^{n \times k}$ where $n = 8$ entities and $k = 10$ functional categories. Each row represents an entity's functional profile vector.

\subsection{Stage 2: Market Tensor Construction}

\subsubsection{Data Collection}

We collect market data via CCXT from Binance:

\begin{itemize}[leftmargin=*]
    \item \textbf{Duration}: One year (October 2024 -- October 2025)
    \item \textbf{Granularity}: Hourly OHLCV candles
    \item \textbf{Assets}: 8 cryptocurrencies matching NLP universe
    \item \textbf{Total observations}: 70,088 rows (8,761 hours $\times$ 8 assets)
\end{itemize}

\subsubsection{Tensor Structure}

We represent market microstructure as a 3-way tensor $\mathcal{X} \in \mathbb{R}^{T \times A \times F}$:

\begin{itemize}[leftmargin=*]
    \item \textbf{Mode 1 (Time)}: $T = 8761$ hourly timestamps
    \item \textbf{Mode 2 (Asset)}: $A = 8$ cryptocurrencies matching NLP universe
    \item \textbf{Mode 3 (Feature)}: $F = 5$ market features (open, high, low, close, volume)
\end{itemize}

The resulting tensor has shape $(8761, 8, 5)$ with 350,440 elements. Z-score normalization is applied per asset-feature pair.

\subsubsection{Normalization Variants}

To ensure robustness across different modeling assumptions, we create three tensor variants:

\begin{enumerate}[leftmargin=*]
    \item \textbf{Raw OHLCV}: No normalization, absolute price levels preserved
    \item \textbf{Z-score normalized}: Per-asset, per-feature standardization $x_{n} = (x - \mu) / \sigma$, enabling cross-asset comparison
    \item \textbf{Log returns}: Stationary time series with features $[\log(p_t / p_{t-1}), \mathrm{high-low range}, \log(V)]$, suitable for returns-based modeling
\end{enumerate}

Our primary analysis uses the z-score normalized tensor for interpretability, with log returns used for robustness checks.

\subsection{Tensor Decomposition}

\subsubsection{CP Decomposition (CANDECOMP/PARAFAC)}

CP decomposition factorizes our 3-way tensor as a sum of rank-one components \citep{Kolda2009Tensor}:

\begin{equation}
\mathcal{X} \approx \sum_{r=1}^{R} \lambda_r \cdot (\mathbf{t}_r \otimes \mathbf{a}_r \otimes \mathbf{f}_r)
\label{eq:cp}
\end{equation}

where $\lambda_r$ are component weights, $\mathbf{t}_r \in \mathbb{R}^T$ (temporal factors), $\mathbf{a}_r \in \mathbb{R}^A$ (asset factors), $\mathbf{f}_r \in \mathbb{R}^F$ (feature factors), and $\otimes$ denotes the outer product.

For this pilot, we construct market profiles from summary statistics (mean return, volatility, Sharpe ratio, max drawdown, average volume, volume volatility, trend) computed from hourly OHLCV data. This yields a market matrix $\mathbf{Y} \in \mathbb{R}^{n \times 7}$ where each row represents an entity's market characteristics. Future work will employ full CP decomposition to extract latent market factors.

\subsubsection{Market Matrix Construction}

Each entity's market profile is computed from one year of hourly Binance data, yielding a market matrix $\mathbf{Y} \in \mathbb{R}^{8 \times 7}$ with standardized features (zero mean, unit variance). This representation captures observable market characteristics that can be compared against whitepaper claims.

\subsection{Stage 3: Alignment Testing}

\subsubsection{Procrustes Rotation}

To compare claims and market spaces with different dimensionalities ($\mathbf{X} \in \mathbb{R}^{8 \times 10}$ vs. $\mathbf{Y} \in \mathbb{R}^{8 \times 7}$), we apply Procrustes rotation \citep{Gower1975Generalized}:

\begin{equation}
\mathbf{R}^* = \arg\min_{\mathbf{R}} \|\mathbf{Y} - \mathbf{X}\mathbf{R}\|_F \quad \text{s.t.} \quad \mathbf{R}^T\mathbf{R} = \mathbf{I}
\label{eq:procrustes}
\end{equation}

The optimal rotation $\mathbf{R}^*$ is computed via SVD of $\mathbf{Y}^T\mathbf{X}$. Both matrices are standardized (zero mean, unit variance per column) before alignment.

\subsubsection{Tucker's Congruence Coefficient}

We measure factor similarity using Tucker's congruence coefficient $\phi$ \citep{Tucker1951Method,Lorenzo-Seva2006Tucker}:

\begin{equation}
\phi = \frac{\sum_i x_i y_i}{\sqrt{\sum_i x_i^2 \cdot \sum_i y_i^2}}
\label{eq:tucker-phi}
\end{equation}

Unlike Pearson correlation, $\phi$ does not subtract means, making it appropriate for comparing factor loadings that may have meaningful zero points.

\textbf{Interpretation thresholds} \citep{Lorenzo-Seva2006Tucker}:
\begin{itemize}[leftmargin=*]
    \item $\phi \geq 0.95$: Factors equivalent
    \item $0.85 \leq \phi < 0.95$: Fair similarity
    \item $0.65 \leq \phi < 0.85$: Some similarity
    \item $\phi < 0.65$: Factors distinct
\end{itemize}

\subsubsection{Bootstrap Confidence Intervals}

We compute 95\% confidence intervals via bootstrap resampling:

\begin{enumerate}[leftmargin=*]
    \item Sample $n$ entities with replacement
    \item Recompute Procrustes alignment and $\phi$ on bootstrap sample
    \item Repeat 50 times
    \item Report 2.5th and 97.5th percentiles
\end{enumerate}

\subsection{Implementation}

All decompositions are implemented using TensorLy v0.8.1 \citep{Kossaifi2019TensorLy} with NumPy 1.26.4 backend. Experiments executed on Resurrexi Lab, a distributed computing cluster with professional Kubernetes orchestration across consumer-grade hardware (8 nodes, 66 cores, 229GB RAM, 48GB VRAM). CP decomposition via ALS with convergence tolerance $10^{-4}$, maximum 500 iterations (typically converges in 50--150 iterations). Tucker decomposition via HOOI with identical tolerance. Random seed set to 42 for reproducibility. CP and Tucker decompositions converge in 30--90 seconds for our dataset.

Code, data, and complete reproducibility instructions will be made publicly available at \url{https://github.com/studiofarzulla/tensor-defi} upon publication.

\section{Results}
\label{sec:results}

\subsection{Functional Profile Extraction}

Zero-shot classification of whitepapers reveals distinct functional profiles for each entity:

\begin{table*}[htbp]
\centering
\small
\caption{Entity Functional Profiles from Whitepaper Classification}
\label{tbl:profiles}
\begin{tabular}{llcc}
\toprule
\textbf{Asset} & \textbf{Project} & \textbf{Primary Function} & \textbf{Top Score} \\
\midrule
BTC & Bitcoin & Medium of Exchange & 45.9\% \\
ETH & Ethereum & Generalist & (distributed) \\
SOL & Solana & Infrastructure & 52.0\% \\
AVAX & Avalanche & Infrastructure & 64.1\% \\
DOT & Polkadot & Governance & 63.3\% \\
FIL & Filecoin & Infrastructure & 50.9\% \\
LINK & Chainlink & Data Oracle & 54.2\% \\
ALGO & Algorand & Identity & 45.4\% \\
\bottomrule
\end{tabular}
\end{table*}

Notable findings: Bitcoin's whitepaper emphasizes ``peer-to-peer electronic cash'' (medium of exchange), not store of value---despite market perception. Ethereum shows no dominant category, reflecting its generalist ``world computer'' positioning. Infrastructure projects (SOL, AVAX, FIL) cluster together with 50-64\% scores.

\begin{figure*}[htbp]
    \centering
    \makebox[\textwidth]{\includegraphics[width=1.15\textwidth]{functional_profiles_heatmap.png}}
    \caption{Heatmap of entity functional profiles from whitepaper zero-shot classification. Darker cells indicate stronger claims in that functional category. Note the infrastructure cluster (SOL/AVAX/FIL), governance emphasis (DOT), and oracle specialization (LINK).}
    \label{fig:profiles-heatmap}
\end{figure*}

\subsection{Market Profile Structure}

Market profiles computed from summary statistics yield a 7-dimensional characterization of each entity's market behavior:

\begin{figure*}[htbp]
    \centering
    \makebox[\textwidth]{\includegraphics[width=1.0\textwidth]{market_profiles.png}}
    \caption{Entity market profiles across 7 summary statistics (mean return, volatility, Sharpe ratio, max drawdown, average volume, volume volatility, trend). Each row represents an entity's standardized position in market characteristic space. Note FIL's extreme positive values across most metrics.}
    \label{fig:market-profiles}
\end{figure*}

\subsection{Alignment Results}

\subsubsection{Overall Congruence}

Procrustes rotation followed by Tucker's congruence yields:

\begin{equation*}
\phi = 0.719 \quad \text{(Some similarity)}
\end{equation*}

\textbf{Bootstrap 95\% CI:} [0.623, 0.953]

The lower bound exceeds the null threshold ($\phi = 0.65$), allowing rejection of H$_0$. Whitepaper functional profiles show statistically significant alignment with market profiles.

\subsubsection{Per-Dimension Congruence}

After Procrustes alignment, we compute Tucker's $\phi$ for each aligned dimension. Substantial heterogeneity exists:

\begin{table}[H]
\centering
\small
\caption{Per-Dimension Tucker's Congruence (Post-Alignment)}
\label{tbl:per-factor}
\begin{tabular}{ccc}
\toprule
\textbf{Dimension} & \textbf{$\phi$} & \textbf{Interpretation} \\
\midrule
Dim 1 & 0.301 & Distinct (claims $\neq$ market) \\
Dim 2 & 0.867 & Fair similarity \\
Dim 3 & 0.728 & Some similarity \\
Dim 4 & 0.982 & Near-equivalent \\
\bottomrule
\end{tabular}
\end{table}

\begin{figure*}[htbp]
    \centering
    \makebox[\textwidth]{\includegraphics[width=1.0\textwidth]{per_dimension_congruence.png}}
    \caption{Per-dimension Tucker's congruence coefficient with interpretation bands. Dimension 4 shows near-perfect alignment ($\phi = 0.98$), while Dimension 1 shows poor alignment ($\phi = 0.30$), suggesting that static founding mandates diverge from evolved market utility in that dimension. Note: with N=8 entities, individual dimension interpretations are illustrative only.}
    \label{fig:per-factor}
\end{figure*}

\subsubsection{Claims vs. Market Space Visualization}

Figure~\ref{fig:claims-market} projects both spaces to 2D via PCA for visualization:

\begin{figure*}[htbp]
    \centering
    \makebox[\textwidth]{\includegraphics[width=1.15\textwidth]{claims_vs_market_space.png}}
    \caption{Claims space (left), market space (center), and Procrustes-aligned overlay (right). Arrows show displacement between claimed and market-perceived positioning. Longer arrows indicate greater narrative-market divergence.}
    \label{fig:claims-market}
\end{figure*}

\section{Discussion}
\label{sec:discussion}

\subsection{Interpreting the Alignment Gap}

The overall $\phi = 0.719$ indicates that markets \textit{partially} price cryptocurrencies according to their stated functional purposes. This ``some similarity'' finding has nuanced implications:

\textbf{Whitepapers are not cheap talk.} The statistically significant alignment (CI excludes $\phi < 0.65$) suggests functional claims contain decision-relevant information that markets incorporate into pricing. Projects claiming similar functions (infrastructure cluster: SOL/AVAX/FIL) exhibit similar market dynamics.

\textbf{But whitepapers are not the whole story.} The $\phi = 0.719$ is far from equivalence ($\phi \geq 0.95$). Other factors---liquidity, speculation, network effects, regulatory exposure---contribute substantially to market positioning independent of stated function.

\subsection{The Per-Dimension Heterogeneity}

The dramatic variation across aligned dimensions ($\phi$ from 0.30 to 0.98) reveals which aspects of claims-market correspondence vary:

\textbf{Dimension 4 ($\phi = 0.982$):} Near-perfect alignment. This aligned dimension shows strong correspondence between whitepaper claims and market characteristics.

\textbf{Dimension 1 ($\phi = 0.301$):} Poor alignment. This dimension represents a gap between static founding mandates and evolved market utility. Possible interpretations:
\begin{itemize}[leftmargin=*]
    \item \textbf{Evolved utility}: Assets developed uses beyond their founding vision (e.g., BTC as ``digital gold'' rather than ``peer-to-peer cash'')
    \item \textbf{Emergent dynamics}: Market characteristics that emerged post-launch (liquidity, exchange listings, regulatory status) are invisible to whitepaper language
    \item \textbf{Speculative dynamics}: Price behavior driven by speculation orthogonal to functional claims
\end{itemize}

\textbf{Caveat:} With N=8 entities, we cannot reliably interpret individual dimension meanings. These observations are illustrative, not inferential.

\subsection{Bitcoin: Illustrating Static Mandate vs. Evolved Utility}

Bitcoin exemplifies the core phenomenon: its whitepaper emphasizes ``peer-to-peer electronic cash'' (medium of exchange: 45.9\%), yet market participants now treat BTC primarily as ``digital gold'' (store of value). This is not ``marketing vs. reality''---Satoshi's vision was genuine---but rather \textit{evolved utility}: the asset's use case shifted over time as the community, market participants, and narrative coalesced around different functions.

This framing matters: we are not detecting \textit{deception} (projects lying in whitepapers) but rather \textit{evolution} (founding mandates becoming outdated as assets mature). The NLP pipeline captures \textit{static founding mandates}, while market profiles capture \textit{current evolved utility}. The alignment coefficient measures the gap between them.

\subsection{Implications}

\subsubsection{Due Diligence}

Projects with large narrative-market gaps warrant scrutiny. Entities showing poor alignment on specific dimensions may be candidates for deeper fundamental analysis---do their claims match market reality?

\subsubsection{Market Efficiency}

The $\phi = 0.719$ suggests semi-strong form efficiency with respect to whitepaper information: markets incorporate functional claims, but incompletely. Systematic trading on narrative-market divergence could exploit this inefficiency.

\subsubsection{Regulatory Classification}

If regulators classify tokens based on stated function, our findings suggest this partially reflects market reality. However, the poor alignment on some dimensions implies certain tokens may be classified based on claims that markets don't validate.

\subsection{Limitations}

\textbf{This section is critical.} The pilot application has several fundamental limitations that preclude strong empirical conclusions. We present them transparently to clarify the paper's scope as methodological, not empirical.

\subsubsection{Sample Size (N=8)}

The most serious limitation. With only 8 entities:
\begin{itemize}[leftmargin=*]
    \item Bootstrap confidence intervals are very wide ([0.62, 0.95])
    \item Individual factor interpretations are unreliable
    \item The $\phi = 0.719$ point estimate carries little inferential weight
    \item Power to detect alignment differences is minimal
\end{itemize}

\textbf{Recommendation for replication:} Future empirical work should use N $\geq$ 30 entities, ideally 50+, for robust alignment inference.

\subsubsection{Feature Collinearity (OHLC)}

Our market tensor uses 5 OHLCV features, but Open, High, Low, and Close are highly collinear within each time period. Effectively, we have only 2-3 independent information dimensions:
\begin{itemize}[leftmargin=*]
    \item \textbf{Price level}: Close (or any OHLC component) captures the main signal
    \item \textbf{Volatility}: High-Low range provides intraperiod dispersion
    \item \textbf{Liquidity}: Volume captures trading activity
\end{itemize}

Future work should replace OHLCV with orthogonalized features: $\{\text{Close}, \log(\text{Volume}), \text{Volatility} = \frac{\text{High}-\text{Low}}{\text{Close}}, \text{Returns} = \log(p_t/p_{t-1})\}$ to avoid inflating factor dimensionality with redundant information.

\subsubsection{Temporal Mismatch}

Whitepapers are static founding documents (written 2008--2021), while market data spans October 2024--October 2025. This creates an inherent asymmetry: we are comparing \textit{founding mandates} with \textit{current market dynamics}.

This is not necessarily a flaw---measuring how well founding visions predict current pricing is precisely the research question---but it should inform interpretation. Low alignment may reflect \textit{evolution} (assets developed new uses) rather than \textit{misrepresentation} (projects lied about their purpose).

\subsubsection{Factor Labeling Circularity}

A methodological concern: if we extract market factors via tensor decomposition and then \textit{interpret} those factors using domain knowledge that overlaps with whitepaper categories (e.g., ``infrastructure factor,'' ``governance factor''), we risk circular reasoning. The current analysis labels factors post-hoc based on entity loadings, but this interpretation is not validated independently.

\textbf{Recommendation:} Future work should identify market factors using exogenous criteria (cross-sectional predictability, macroeconomic sensitivity) \textit{before} alignment testing, to avoid interpreting market factors in terms of whitepaper categories.

\subsubsection{Additional Limitations}

\begin{itemize}[leftmargin=*]
    \item \textbf{Zero-shot accuracy}: BART-MNLI may miss crypto-specific nuance; fine-tuned models could improve classification
    \item \textbf{Single venue}: Binance only; cross-exchange dynamics may differ
    \item \textbf{Whitepaper selection}: Only projects with available whitepapers; excludes memecoins, newer projects
\end{itemize}

\subsection{Future Work}

Given these limitations, future work will focus on:

\begin{enumerate}[leftmargin=*]
    \item \textbf{Expanded universe}: 50+ entities for robust empirical inference
    \item \textbf{Feature engineering}: Orthogonalized market features (Close, Volume, Volatility, Returns) instead of collinear OHLC
    \item \textbf{Independent factor identification}: Identify market factors via exogenous criteria before alignment testing
    \item \textbf{Temporal analysis}: Rolling window alignment to detect how narrative-market gaps evolve
    \item \textbf{Causal analysis}: Granger causality tests for whether claims lead or lag market positioning
    \item \textbf{Fine-tuned NLP}: Domain-specific models trained on crypto documentation (CryptoBERT)
\end{enumerate}

\section{Conclusion}
\label{sec:conclusion}

We have proposed a methodological framework for testing whether cryptocurrency whitepaper claims align with market factor structures. By combining NLP-based functional taxonomy extraction with tensor decomposition of market dynamics, we measure the correspondence between what projects \textit{claim} in their founding documents and how markets \textit{currently price} them.

\textbf{Methodological Contribution:} The three-stage framework (NLP extraction $\rightarrow$ tensor decomposition $\rightarrow$ Procrustes/Tucker alignment) provides a principled approach to testing narrative-market correspondence. This is the paper's primary contribution.

\textbf{Illustrative Application:} The pilot with N=8 entities yields $\phi = 0.719$ (95\% CI: [0.62, 0.95]), but the small sample and wide CI preclude strong empirical claims. The point estimate is illustrative, not inferential.

\textbf{Key Conceptual Framing:} We interpret divergence as \textit{static mandate vs. evolved utility}, not ``marketing vs. reality.'' Bitcoin's whitepaper describes ``peer-to-peer electronic cash,'' but the asset's use evolved toward ``digital gold.'' The gap represents \textit{evolution}, not deception. This framing is essential for correctly interpreting alignment results.

\textbf{Limitations Acknowledged:} The pilot has fundamental limitations: N=8 sample size, OHLC collinearity (effectively 2-3 features, not 5), temporal mismatch between static whitepapers and current markets, and potential circularity in factor labeling. These are documented for future replication efforts.

\textbf{Future Empirical Work:} Robust empirical conclusions require N $\geq$ 30 entities, orthogonalized features, and independent factor identification. The methodology established here provides the foundation; the empirical work remains to be done.

% ====================
% SWITCH BACK TO SINGLE-COLUMN FOR REFERENCES
% ====================
\clearpage
\onecolumn
\setstretch{1.5}  % Re-apply line spacing

% ====================
% DATA AVAILABILITY & ACKNOWLEDGEMENTS
% ====================

\section*{Data Availability Statement}

All data, analysis code, figure generation scripts, and comprehensive documentation are publicly available at: \url{https://github.com/studiofarzulla/tensor-defi}. The repository includes:

\begin{itemize}[leftmargin=*]
    \item Whitepaper PDFs and extracted text
    \item NLP pipeline code (zero-shot classification)
    \item Tensor decomposition and alignment testing code
    \item Pre-computed functional profiles and alignment results
    \item Figure generation scripts
\end{itemize}

Data sources include cryptocurrency whitepapers (public documents) and Binance historical OHLCV data (via CCXT library). Models are implemented in Python 3.13+ using TensorLy, Transformers, and SciPy.

\section*{Acknowledgements}

This research benefited from Perplexity AI for efficient literature discovery and source verification, and Anthropic's Claude for assistance with analytical framework development, custom tensor decomposition implementation, literature synthesis, and technical writing. All conceptual innovations, theoretical claims, and interpretive judgments remain the author's sole responsibility.

All computational analysis was conducted at Resurrexi Lab, a distributed computing cluster built from consumer-grade hardware (8 nodes, 66 cores, 229GB RAM, 48GB VRAM), demonstrating that rigorous quantitative finance research is accessible without institutional supercomputing infrastructure.

Research methodologies and reproducibility practices are documented at \url{https://farzulla.org/methodologies}.

% ====================
% BIBLIOGRAPHY
% ====================
\bibliography{references}

\end{document}
